\section{Introducao}

A gestao financeira pessoal e um desafio enfrentado por milhoes de pessoas ao redor do mundo. A falta de organizacao e controle sobre receitas e despesas pode levar a problemas financeiros graves, endividamento e dificuldades para alcancar objetivos de vida.

\subsection{Contexto do Problema}

Atualmente, muitas pessoas organizam suas financas utilizando **planilhas Excel** ou metodos manuais, que podem ser ineficientes, propensos a erros e dificeis de acessar em dispositivos moveis. O FinOps surge como uma **evolucao digital dessas planilhas**, mantendo o controle manual que os usuarios conhecem, mas oferecendo melhor organizacao, visualizacao e acessibilidade.

A necessidade de uma ferramenta que **substitua o Excel** para gestao financeira se torna evidente, especialmente considerando:

\begin{itemize}
    \item A dificuldade de usar planilhas complexas em dispositivos moveis
    \item A necessidade de colaboracao financeira entre casais e grupos
    \item A importancia de visualizacoes claras e automaticas dos dados
    \item A demanda por uma interface intuitiva sem complexidade excessiva
    \item A necessidade de **nao depender de conexoes bancarias** ou APIs externas
\end{itemize}

\subsection{Motivacao}

O projeto FinOps nasce da necessidade de **digitalizar e modernizar as planilhas Excel** de controle financeiro, criando uma solucao pratica e intuitiva. O sistema funciona com base em **indicacoes percentuais** e entrada manual de dados, seguindo um fluxo logico e eficiente:

\begin{enumerate}
    \item Registro manual de todas as receitas mensais
    \item Deducao dos gastos fixos obrigatorios
    \item Alocacao percentual para custos variaveis
    \item **Indicacao percentual** para poupanca e reserva de emergencia
    \item **Sugestao percentual** para investimentos (sem acompanhamento automatico)
    \item Calculo e divisao dos gastos livres restantes
\end{enumerate}

\textbf{Importante}: O sistema **nao se conecta** a bancos, corretoras ou APIs externas. Funciona como um **organizador digital** onde o usuario insere manualmente os dados, similar ao Excel, mas com interface moderna e acessivel.

\subsection{Justificativa}

A criacao de uma aplicacao para **substituir planilhas Excel** na gestao financeira se justifica pela:

\begin{itemize}
    \item \textbf{Modernizacao}: Interface moderna e responsiva que funciona bem em qualquer dispositivo
    \item \textbf{Simplicidade}: Mantem a simplicidade do Excel, mas com melhor organizacao visual
    \item \textbf{Acessibilidade}: Disponivel em web e mobile para uso a qualquer momento
    \item \textbf{Colaboracao}: Suporte para gestao financeira compartilhada (casais, republicas)
    \item \textbf{Visualizacao}: Graficos automaticos e relatorios para melhor compreensao dos dados
    \item \textbf{Independencia}: Nao depende de conexoes bancarias ou APIs externas
    \item \textbf{Controle Total}: O usuario mantem controle total sobre seus dados e entrada manual
\end{itemize}
